\documentclass[aspectratio=169]{beamer}

\usepackage[utf8]{inputenc}
\usepackage[T1]{fontenc}
\usepackage{amsmath,amssymb,amsfonts,bm,mathtools}
\usepackage{physics}
\usepackage{braket}
\usepackage{hyperref}

\usetheme{Madrid}

\title{CTNN Backflow vs Standard Backflow}
\subtitle{Modification of the Slater Determinant and Connection to the Variational Principle}
\author{}
\date{}

\begin{document}

\frame{\titlepage}

\begin{frame}{Aim of These Notes}
These slides explain:
\begin{enumerate}
    \item what standard backflow is,
    \item how CTNN backflow differs from standard backflow,
    \item how backflow modifies a Slater determinant,
    \item and how the full construction enters the variational principle.
\end{enumerate}

The main physical context is a fermionic many-body wave function, where one wants a flexible yet antisymmetric ansatz that captures strong correlations and improves the nodal surface.
\end{frame}

\section{Standard backflow}

\begin{frame}{Ordinary Slater-Jastrow Ansatz}
A common fermionic variational ansatz is
\[
\Psi(\mathbf{R}) = e^{J(\mathbf{R})}\det[\phi_a(\mathbf{r}_i)],
\qquad
\mathbf{R}=(\mathbf{r}_1,\dots,\mathbf{r}_N).
\]

Here:
\begin{itemize}
    \item \(J(\mathbf{R})\) is a symmetric Jastrow factor,
    \item the determinant enforces antisymmetry,
    \item the orbitals are evaluated at the bare coordinates \(\mathbf{r}_i\).
\end{itemize}

Backflow modifies these coordinates before they enter the determinant.
\end{frame}

\begin{frame}{Standard Backflow}
Standard backflow introduces effective coordinates
\[
\tilde{\mathbf{r}}_i = \mathbf{r}_i + \bm{\xi}_i(\mathbf{R}),
\]
with a typical ansatz
\[
\bm{\xi}_i(\mathbf{R})
=
\sum_{j\neq i} \eta(r_{ij})(\mathbf{r}_i-\mathbf{r}_j),
\qquad
r_{ij}=|\mathbf{r}_i-\mathbf{r}_j|.
\]

So particle \(i\) is displaced by a sum of pairwise contributions generated by the other particles.
\end{frame}

\begin{frame}{Properties of Standard Backflow}
Standard backflow is:
\begin{itemize}
    \item physically motivated,
    \item usually pairwise,
    \item analytically structured,
    \item and relatively easy to interpret.
\end{itemize}

Its limitations are:
\begin{itemize}
    \item restricted functional flexibility,
    \item difficulty capturing higher-body structure,
    \item dependence on a chosen analytic form.
\end{itemize}
\end{frame}

\section{CTNN backflow}

\begin{frame}{CTNN Backflow: Core Idea}
CTNN backflow replaces the fixed analytic pairwise form by a learned many-body map:
\[
\tilde{\mathbf{r}}_i = \mathbf{r}_i + F_{i,\theta}(\mathbf{R}),
\]
or collectively,
\[
\tilde{\mathbf{R}} = F_\theta(\mathbf{R}).
\]

So the essential change is
\[
\text{analytic pairwise deformation}
\quad\longrightarrow\quad
\text{learned many-body coordinate transformation}.
\]
\end{frame}

\begin{frame}{How CTNN Backflow Differs}
A useful comparison is:

\begin{center}
\begin{tabular}{lcc}
\hline
Feature & Standard backflow & CTNN backflow \\
\hline
Functional form & analytic & learned \\
Typical structure & pairwise & many-body / hierarchical \\
Parameters & few & many \\
Expressivity & limited & high \\
Interpretability & high & moderate \\
\hline
\end{tabular}
\end{center}

Conceptually:
\[
\boxed{\text{Standard backflow}=\text{analytic pairwise correction}}
\]
\[
\boxed{\text{CTNN backflow}=\text{learned many-body deformation}}
\]
\end{frame}

\begin{frame}{A Generic CTNN Backflow Decomposition}
One may imagine a neural construction of the form
\[
h_i = \phi_\theta(\mathbf{r}_i),\qquad
h_{ij}=\psi_\theta(\mathbf{r}_i,\mathbf{r}_j),
\]
followed by aggregation,
\[
H_i = \sum_{j\neq i} h_{ij},
\]
and finally
\[
\tilde{\mathbf{r}}_i = \mathbf{r}_i + g_\theta(h_i,H_i).
\]

This already goes beyond standard backflow because the functions are learned and can represent nontrivial collective structure.
\end{frame}

\section{Modification of the Slater determinant}

\begin{frame}{Ordinary Slater Determinant}
Without backflow, the Slater matrix is
\[
M_{ia}(\mathbf{R})=\phi_a(\mathbf{r}_i),
\]
and the determinant is
\[
D(\mathbf{R})=\det M(\mathbf{R}).
\]

So each row depends only on the coordinate of one particle.
\end{frame}

\begin{frame}{Backflow-Modified Slater Determinant}
With backflow,
\[
\mathbf{r}_i \to \tilde{\mathbf{r}}_i(\mathbf{R}),
\]
and therefore
\[
\widetilde{M}_{ia}(\mathbf{R})=\phi_a(\tilde{\mathbf{r}}_i(\mathbf{R})),
\]
so the determinant becomes
\[
D_{\mathrm{BF}}(\mathbf{R})
=
\det \widetilde{M}(\mathbf{R})
=
\det[\phi_a(\tilde{\mathbf{r}}_i(\mathbf{R}))].
\]

Now each row depends on the full many-body configuration through \(\tilde{\mathbf{r}}_i(\mathbf{R})\).
\end{frame}

\begin{frame}{What Has Changed Mathematically?}
Originally,
\[
M_{ia}=\phi_a(\mathbf{r}_i),
\]
so the dependence is one-body at the level of the orbital arguments.

After backflow,
\[
\widetilde{M}_{ia}
=
\phi_a\bigl(\mathbf{r}_i + F_{i,\theta}(\mathbf{R})\bigr),
\]
so every row becomes a many-body object.

Thus the determinant is still antisymmetric, but it is now built from configuration-dependent quasi-orbitals rather than bare one-particle orbitals.
\end{frame}

\begin{frame}{Why Antisymmetry is Preserved}
If the backflow map is permutation-equivariant, then exchanging particles \(i\leftrightarrow j\) exchanges the corresponding effective coordinates:
\[
\tilde{\mathbf{r}}_i(\mathbf{R}) \leftrightarrow \tilde{\mathbf{r}}_j(\mathbf{R}).
\]

Hence two rows of the determinant are exchanged, giving
\[
D_{\mathrm{BF}}(\mathbf{R})\mapsto -D_{\mathrm{BF}}(\mathbf{R}).
\]

So antisymmetry is preserved.
\end{frame}

\begin{frame}{Backflow Plus Jastrow}
A common full ansatz is
\[
\Psi_\theta(\mathbf{R})
=
e^{J_\theta(\mathbf{R})}
\det[\phi_a(\tilde{\mathbf{r}}_i(\mathbf{R}))].
\]

Interpretation:
\begin{itemize}
    \item \(J_\theta\) handles symmetric correlations,
    \item the determinant enforces antisymmetry,
    \item backflow changes the nodal structure through the effective coordinates.
\end{itemize}

This is important because Jastrow factors alone do not move the nodes, while backflow does.
\end{frame}

\begin{frame}{Small-Displacement Expansion}
If
\[
\tilde{\mathbf{r}}_i=\mathbf{r}_i+\delta \mathbf{r}_i,
\]
with a small backflow displacement, then
\[
\phi_a(\tilde{\mathbf{r}}_i)
\approx
\phi_a(\mathbf{r}_i)
+
\delta \mathbf{r}_i \cdot \nabla \phi_a(\mathbf{r}_i)
+\cdots.
\]

Thus
\[
\widetilde{M}_{ia}
\approx
\phi_a(\mathbf{r}_i)
+
F_{i,\theta}(\mathbf{R})\cdot \nabla\phi_a(\mathbf{r}_i).
\]

So backflow introduces configuration-dependent gradient corrections inside the determinant.
\end{frame}

\section{Connection to the variational principle}

\begin{frame}{Variational Principle}
Given a Hamiltonian \(H\), the variational energy of a trial state \(\Psi_\theta\) is
\[
E(\theta)
=
\frac{\bra{\Psi_\theta}H\ket{\Psi_\theta}}{\braket{\Psi_\theta|\Psi_\theta}}.
\]

The variational principle states
\[
E(\theta)\ge E_0,
\]
where \(E_0\) is the exact ground-state energy.

So the goal is to minimize \(E(\theta)\) over the ansatz parameters.
\end{frame}

\begin{frame}{Backflow as Variational Degrees of Freedom}
With CTNN backflow, the parameters \(\theta\) include those defining
\[
\tilde{\mathbf{r}}_i(\mathbf{R})=\mathbf{r}_i+F_{i,\theta}(\mathbf{R}).
\]

So the variational manifold is enlarged beyond:
\begin{itemize}
    \item Jastrow parameters,
    \item orbital coefficients,
\end{itemize}
to include:
\begin{itemize}
    \item learned many-body geometry of the effective coordinates.
\end{itemize}

This additional flexibility can lower the energy.
\end{frame}

\begin{frame}{Why Backflow Can Improve the Variational State}
Backflow helps in two main ways:
\begin{enumerate}
    \item it captures additional many-body correlations,
    \item it changes the nodal surface of the fermionic wave function.
\end{enumerate}

The nodal surface is especially important in fermionic systems, since it strongly influences the quality of the wave function and the achievable ground-state energy.
\end{frame}

\begin{frame}{Local Energy}
In variational Monte Carlo one uses the local energy
\[
E_{\mathrm{loc}}(\mathbf{R})
=
\frac{H\Psi_\theta(\mathbf{R})}{\Psi_\theta(\mathbf{R})}.
\]

Then
\[
E(\theta)
=
\frac{\int d\mathbf{R}\, |\Psi_\theta(\mathbf{R})|^2 E_{\mathrm{loc}}(\mathbf{R})}
{\int d\mathbf{R}\, |\Psi_\theta(\mathbf{R})|^2}.
\]

Thus optimization requires evaluating the backflow-modified determinant and its derivatives on sampled configurations.
\end{frame}

\begin{frame}{Derivative Structure}
Because
\[
\widetilde{M}_{ia}=\phi_a(\tilde{\mathbf{r}}_i(\mathbf{R})),
\]
spatial derivatives involve the chain rule:
\[
\nabla_{\mathbf{r}_k}\widetilde{M}_{ia}
=
\sum_{\alpha}
\frac{\partial \phi_a}{\partial \tilde{r}_{i,\alpha}}
\frac{\partial \tilde{r}_{i,\alpha}}{\partial \mathbf{r}_k}.
\]

Therefore CTNN backflow modifies:
\begin{itemize}
    \item gradients,
    \item Laplacians,
    \item kinetic-energy contributions,
    \item and parameter gradients for optimization.
\end{itemize}
\end{frame}

\begin{frame}{Variational Optimization Problem}
The full optimization problem is
\[
\theta^\star
=
\arg\min_\theta
\frac{\bra{\Psi_\theta}H\ket{\Psi_\theta}}{\braket{\Psi_\theta|\Psi_\theta}},
\]
with
\[
\Psi_\theta(\mathbf{R})
=
e^{J_\theta(\mathbf{R})}
\det[\phi_a(\mathbf{r}_i+F_{i,\theta}(\mathbf{R}))].
\]

So CTNN backflow is not merely a preprocessing step: it is an intrinsic part of the variational wave function.
\end{frame}

\section{Summary}

\begin{frame}{Summary of the Main Differences}
\begin{itemize}
    \item Standard backflow uses analytic pairwise coordinate corrections.
    \item CTNN backflow uses a learned many-body map \(F_{i,\theta}(\mathbf{R})\).
    \item The Slater determinant is modified by evaluating orbitals at effective coordinates:
    \[
    \phi_a(\mathbf{r}_i)\to \phi_a(\tilde{\mathbf{r}}_i(\mathbf{R})).
    \]
    \item Antisymmetry is preserved if the map is permutation-equivariant.
    \item In the variational principle, CTNN backflow enlarges the manifold of trial states and can lower the ground-state energy by improving many-body correlations and nodal structure.
\end{itemize}
\end{frame}

\begin{frame}{Final Take-Home Message}
The essential formulas are
\[
\tilde{\mathbf{r}}_i=\mathbf{r}_i+F_{i,\theta}(\mathbf{R}),
\]
\[
\Psi_\theta(\mathbf{R})
=
e^{J_\theta(\mathbf{R})}
\det[\phi_a(\tilde{\mathbf{r}}_i(\mathbf{R}))],
\]
\[
E(\theta)=\frac{\bra{\Psi_\theta}H\ket{\Psi_\theta}}{\braket{\Psi_\theta|\Psi_\theta}}.
\]

These summarize CTNN backflow as a learned many-body coordinate transformation embedded directly inside an antisymmetric variational ansatz.
\end{frame}

\end{document}
